\begin{textsample}{POS Dim 2 – 2017\_c4 – Score 11.00 – 2017\_c4/enade\_2017\_ci\_bio\_3400d20d88.txt}  \label{ex:f2_pos_096}
CAPÍTULO 4 PERCEPÇÃO DA PROVA As análises feitas neste capítulo tratam das percepções dos concluintes da Área de Ciências Biológicas sobre a prova aplicada no Enade/2017.
Essas percepções foram mensuradas por meio de nove questões que avaliaram desde o grau de dificuldade da prova até o tempo gasto para concluí -la.
As percepções sobre a prova foram relacionadas ao desempenho dos estudantes e à Grande Região de funcionamento do curso.
O \textbf{questionário} de percepção da prova encontra -se ao final do Anexo VII, que traz a \textbf{reprodução} do exame.
O desempenho dos estudantes foi classificado em quatro quartos.
Para tanto, esse desempenho foi \textbf{ordenado} de forma ascendente.
O percentil 25, P25, também conhecido como primeiro quartil, é a nota de desempenho que deixa um quarto ( 25 \% ) dos valores observados abaixo e três quartos acima.
A Figura 1 apresenta uma ilustração desse conceito.
O quarto inferior de desempenho é composto pelas notas abaixo do primeiro quartil.
Já o percentil 75, P75, também conhecido como terceiro quartil, é o valor para o qual há três quartos ( 75 \% ) dos dados abaixo e um quarto acima dele.
O quarto superior de desempenho é composto pelas notas iguais ou acima do terceiro quartil.
O percentil 50, P50, também conhecido como mediana, é o valor que divide as notas em dois conjuntos de igual tamanho.
O segundo quarto inclui valores entre o primeiro quartil ( P25 ) e a mediana.
O terceiro quarto contém os valores entre a mediana ( P50 ) e o terceiro quartil ( P75 ).
Vale ressaltar que percentis, quartis e medianas são pontos que não obrigatoriamente pertencem ao conjunto original de dados, ao passo que os quartos são \textbf{subconjuntos} dos dados originais.
A seguir, serão apresentados gráficos com resultados selecionados, relativos às nove questões avaliadas por grupos de estudantes.
Os gráficos apresentam nas barras o percentual de alunos que assinalaram uma das opções ou a soma das porcentagens daqueles que assinalaram duas ( ou três ) delas.
Por exemplo, para as questões 1 e 2, os gráficos apresentam a porcentagem total de participantes que assinalaram as opções ( D ) Difícil e ( E ) Muito difícil.
Em cada barra foram assinalados também os extremos do intervalo de confiança de 95 \% como linhas verticais unidas por uma linha horizontal na forma da letra H maiúscula.
As Tabelas no Anexo II apresentam os valores absolutos e a distribuição percentual17 das alternativas válidas das nove questões, segundo o mesmo \textbf{recorte} de desempenho dos alunos e Grande Região de funcionamento do curso. 4.1 GRAU DE DIFICULDADE DA PROVA 4.1.1 Componente de Formação Geral Ao avaliarem Qual o grau de dificuldade desta prova na parte de Formação Geral? ( Questão 1 ), 31,3 \% do grupo de inscritos e presentes de Licenciatura optaram pelas alternativas Difícil ou Muito difícil.
Entretanto, para mais da metade dos estudantes ( 59,9 \% ), o Componente de Formação Geral da prova foi considerado com grau de dificuldade Médio ( Gráfico 4.1a, Gráfico 4.2a e, no Anexo II, a Tabela II.1a ).
O percentual de estudantes de Licenciatura que consideraram a prova como Difícil ou Muito difícil foi maior na região Sul, onde a proporção foi de 36,0 \%, enquanto a de menor incidência foi na Sudeste, com 27,9 \%.
No Gráfico 4.1a, é possível observar que as diferenças entre a região Sudeste e as demais regiões são estatisticamente significativas, bem como a diferença entre as regiões Sul e Nordeste.
Nas Grandes Regiões, a proporção de presentes à prova que consideraram o Componente de Formação Geral como sendo de grau de dificuldade Médio esteve entre 55,1 \% na região Centro-Oeste, e 61,5 \% na região Sudeste. 17 Cumpre lembrar uma das convenções para tabelas \textbf{numéricas} ( pág. iii ) sobre a possibilidade de a soma das partes não resultar em 100 \%, por questões de \textbf{arredondamento}.
Ao avaliarem Qual o grau de dificuldade desta prova na parte de Formação Geral? ( Questão 1 ), 19,6 \% do grupo de inscritos e presentes de Bacharelado optaram pelas alternativas Difícil ou Muito difícil.
Entretanto, para cerca de dois terços dos estudantes ( 66,0 \% ), o Componente de Formação Geral da prova foi considerado com grau de dificuldade Médio ( Gráfico 4.1b, Gráfico 4.2b e, no Anexo II, a Tabela II.1b ).
O percentual de estudantes de Bacharelado que consideraram a prova como Difícil ou Muito difícil foi maior na região Norte, onde a proporção foi de 24,7 \%, enquanto a de menor incidência foi na Nordeste, com 16,8 \%.
No Gráfico 4.1b, é possível observar que as diferenças da proporção dessa resposta na região Nordeste e em relação às regiões Norte e Sul são estatisticamente significativas.
Nas Grandes Regiões, a proporção de presentes à prova que consideraram o Componente de Formação Geral como sendo de grau de dificuldade Médio esteve entre 63,3 \% na região Norte e 67,4 \% na região Nordeste.
O percentual de alunos de Licenciatura que consideraram a prova Difícil ou Muito difícil foi decrescente em função dos quartos de desempenho : 40,9 \% no primeiro quarto e 21,6 \% no quarto, grupo de melhor desempenho na prova.
Nos quartos de desempenho intermediários, a proporção de alunos que consideraram a prova Difícil ou Muito difícil foi de 34,2 \% no segundo quarto e de 29,1 \% no terceiro quarto.
As diferenças entre os quartos são estatisticamente significativas.
Para todos os quartos de desempenho, a alternativa modal para essa pergunta foi Médio, com 52,1 \% e 66,5 \% dos respondentes nos quartos extremos, primeiro e quarto, respectivamente.
O percentual de alunos de Bacharelado que consideraram a prova Difícil ou Muito difícil foi decrescente em função dos quartos de desempenho : 26,1 \% no primeiro quarto e 14,8 \% no quarto, grupo de melhor desempenho na prova.
Nos quartos de desempenho intermediários, a proporção de alunos que consideraram a prova Difícil ou Muito difícil foi de 20,5 \% no segundo quarto e 17,2 \% no terceiro quarto.
As diferenças entre o quarto inferior de desempenho e os demais são estatisticamente significativas, bem como entre o segundo quarto e o quarto superior de desempenho.
Para todos os quartos de desempenho, a alternativa modal para essa pergunta foi Médio, com 62,5 \% e 66,9 \% dos respondentes nos quartos extremos, primeiro e quarto, respectivamente. 4.1.2 Componente de Conhecimento Específico Ao responderem à Questão 2 – Qual o grau de dificuldade desta prova na parte de Componente Específico? – 41,9 \% do grupo de estudantes de Licenciatura classificaram -na como Difícil ou Muito difícil.
Além disso, o Componente de Conhecimento Específico da prova foi considerado com grau de dificuldade Médio por 53,2 \% dos alunos ( Gráfico 4.3a, Gráfico 4.4a, e, no Anexo II, a Tabela II.2a ).
A análise das respostas dos estudantes de Licenciatura quanto ao grau de dificuldade do Componente de Conhecimento Específico da prova, agregado por Grande Região, mostra que a diferença entre a maior e as demais proporções de alunos que a avaliaram como Difícil ou Muito difícil é estatisticamente significativa : a maior na região Sul ( 48,4 \% ) e as demais proporções nas regiões Nordeste ( 42,6 \% ), Norte ( 42,2 \% ), Centro-Oeste ( 41,4 \% ) e Sudeste ( 38,4 \% ).
A A diferença das proporções dessa resposta nas regiões Sudeste e Nordeste também é significativa.
O percentual de alunos que classificaram o grau de dificuldade como Médio, no Componente de Conhecimento Específico, variou de 47,4 \% a 56,4 \%, para as regiões Sul e Sudeste, respectivamente.
Ao responderem à Questão 2 – Qual o grau de dificuldade desta prova na parte de Componente Específico? – 33,8 \% do grupo de estudantes de Bacharelado classificaram -na como Difícil ou Muito difícil.
Além disso, o Componente de Conhecimento Específico da prova foi considerado com grau de dificuldade Médio por 59,1 \% dos alunos ( Gráfico 4.3b, Gráfico 4.4b, e, no Anexo II, a Tabela II.2b ).
A análise das respostas dos estudantes de Bacharelado quanto ao grau de dificuldade do Componente de Conhecimento Específico da prova, agregado por Grande Região, mostra que a proporção de alunos que a avaliaram como Difícil ou Muito difícil na região Nordeste ( 28,9 \% ) quando comparada às proporções das regiões Sudeste ( 33,9 \% ), Sul ( 36,1 \% ) e Centro-Oeste ( 40,1 \% ) apresenta diferença estatisticamente significativa.
O percentual de alunos que classificaram o grau de dificuldade como Médio, no Componente de Conhecimento Específico, variou de 52,0 \% a 62,4 \%, para as regiões Centro-Oeste e Nordeste, respectivamente.
Considerando -se a avaliação da dificuldade das questões do Componente de Conhecimento Específico da prova, de acordo com o desempenho dos estudantes de Licenciatura, observa -se que há diferença estatisticamente significativa dos resultados entre o último e os demais quartos de desempenho, e, também, entre o primeiro e o terceiro quarto.
A proporção dos que classificaram a parte específica como Difícil ou Muito difícil variou de 35,3 \% ( último quarto ) a 46,6 \% ( primeiro quarto ).
A alternativa modal para a Questão 2 foi o grau Médio, com 48,3 \% do quarto inferior e 58,6 \% do superior optando por essa resposta.
Considerando -se a avaliação da dificuldade das questões do Componente de Conhecimento Específico da prova, de acordo com o desempenho dos estudantes de Bacharelado, observa -se que há diferença estatisticamente significativa dos resultados entre o último e os demais quartos de desempenho, bem como entre o primeiro e o terceiro.
A proporção dos que classificaram a parte específica como Difícil ou Muito difícil variou de 27,1 \% ( último quarto ) a 39,5 \% ( primeiro quarto ).
A alternativa modal para a Questão 2 foi o grau Médio, com 54,2 \% do quarto inferior e 62,5 \% do superior optando por essa resposta. 4.2 EXTENSÃO DA PROVA EM RELAÇÃO AO TEMPO TOTAL \textbf{Indagados} quanto à extensão da prova, em relação ao tempo total oferecido para a sua resolução ( Questão 3 ), os estudantes de Licenciatura apontaram, com maior incidência, a alternativa que considerava a extensão adequada, para todas as agregações consideradas ( Gráfico 4.5a, Gráfico 4.6a, e, no Anexo II, a Tabela II.3a ).
O percentual de alunos de Licenciatura que responderam ser a extensão da prova adequada foi de 49,4 \%.
Já 43,5 \% dos inscritos presentes consideraram que a prova foi longa ou muito longa, e 7,1 \% a avaliaram como curta ou muito curta.
Dentre as Grandes Regiões, a proporção daqueles que avaliaram a prova como longa ou muito longa em relação ao tempo total destinado à sua resolução variou de 34,9 \% na região Norte até 47,3 \% na região Sul.
A diferença entre a região Norte e as demais é estatisticamente significativa, e também o são as diferenças das proporções das regiões Nordeste e Centro-Oeste em relação às regiões Sudeste e ao Sul. \textbf{Indagados} quanto à extensão da prova, em relação ao tempo total oferecido para a sua resolução ( Questão 3 ), os estudantes de Bacharelado apontaram, com maior incidência, a alternativa que considerava a extensão adequada, para todas as agregações consideradas ( Gráfico 4.5b, Gráfico 4.6b, e, no Anexo II, a Tabela II.3b ).
O percentual de alunos de Bacharelado que responderam ser a extensão da prova adequada foi de 55,0 \%.
Já 38,9 \% dos inscritos presentes consideraram que a prova foi longa ou muito longa, e 6,1 \% a avaliaram como curta ou muito curta.
Dentre as Grandes Regiões, a proporção daqueles que avaliaram a prova como longa ou muito longa em relação ao tempo total destinado à sua resolução variou de 35,5 \% na região Centro-Oeste até 41,9 \% na região Sul.
Não são observadas diferenças estatisticamente significativas dentre as regiões.
Considerando -se o desempenho dos alunos de Licenciatura, nota -se ainda que 46,9 \% consideraram a extensão da prova adequada no quarto de desempenho inferior e 50,4 \% no de melhor desempenho ( quarto superior ).
Nos quartos intermediários, essa proporção foi de 49,5 \% no segundo quarto e de 50,6 \% no terceiro.
No Gráfico 4.6a, pode -se constatar que não há diferença estatisticamente significativa entre as proporções de estudantes de Licenciatura que consideraram a prova longa ou muito longa em função dos quartos de desempenho.
Considerando -se o desempenho dos alunos de Bacharelado, nota -se ainda que 52,8 \% consideraram a extensão da prova adequada no quarto de desempenho inferior e 55,9 \% no de melhor desempenho ( quarto superior ).
Nos quartos intermediários, essa proporção foi de 54,9 \% no segundo quarto e de 56,5 \% no terceiro.
No Gráfico 4.6b, pode -se constatar que há uma tendência decrescente da proporção de estudantes de Bacharelado que consideraram a prova longa ou muito longa em função dos quartos de desempenho, variando de 40,4 \%, no quarto inferior, a 37,6 \%, no quarto superior.
Observa -se, porém, que não há diferença estatisticamente significativa entre as proporções de escolha desta resposta. 4.3 COMPREENSÃO DOS ENUNCIADOS DAS QUESTÕES 4.3.1 Componente de Formação Geral Em relação aos enunciados das questões do Componente de Formação Geral ( Questão 4 ), as opiniões foram positivas, já que 73,9 \% dos alunos de Licenciatura avaliados consideraram todos ou a maioria dos enunciados das questões claros e objetivos ( Gráfico 4.7a, Gráfico 4.8a, e, no Anexo II, a Tabela II.4a ).
Na análise regional, a \textbf{percentagem} de estudantes de Licenciatura que avaliaram que todos ou a maioria dos enunciados das questões do Componente de Formação Geral estavam claros e objetivos variou de 71,2 \% na região Norte a 75,3 \% nas regiões Sudeste e Sul, observando -se diferença estatisticamente significativa entre a região Sudeste e as regiões Norte e Nordeste.
A análise das percepções dos estudantes de Licenciatura sobre a clareza e objetividade dos enunciados permite afirmar que todos, ou a maioria dos enunciados de questões relativas ao Componente de Formação Geral, foram considerados claros e objetivos para a maior parte dos respondentes ( maior ou igual a 71,2 \% em todas as regiões e maior ou igual a 66,0 \% para todos os quartos de desempenho ).
Em relação aos enunciados das questões do Componente de Formação Geral ( Questão 4 ), as opiniões foram positivas, já que 79,5 \% dos alunos de Bacharelado avaliados consideraram todos ou a maioria dos enunciados das questões claros e objetivos ( Gráfico 4.7b, Gráfico 4.8b, e, no Anexo II, a Tabela II.4b ).
Na análise regional, a \textbf{percentagem} de estudantes de Bacharelado que avaliaram que todos ou a maioria dos enunciados das questões do Componente de Formação Geral estavam claros e objetivos variou de 77,9 \% na região Sul a 82,0 \% na região Norte, não havendo diferenças estatisticamente significativas entre as regiões.
A análise das percepções dos estudantes de Bacharelado sobre a clareza e objetividade dos enunciados permite afirmar que todos, ou a maioria dos enunciados de questões relativas ao Componente de Formação Geral, foram considerados claros e objetivos para a maior parte dos respondentes ( maior ou igual a 77,9 \% em todas as regiões e maior ou igual a 74,6 \% para todos os quartos de desempenho ).
Segundo o desempenho, observa -se que a proporção dos alunos de Licenciatura que emitiram essa opinião \textbf{cresce} conforme o desempenho aumenta, com diferenças estatisticamente significativas entre todos os quartos de desempenho.
No quarto superior, a clareza e objetividade de todos ou da maioria dos enunciados das questões foi percebida por 79,9 \%.
Segundo o desempenho, observa -se que a proporção dos alunos de Bacharelado que emitiram essa opinião \textbf{cresce} conforme o desempenho aumenta, com diferenças estatisticamente significativas entre todos os quartos de desempenho.
No quarto superior, a clareza e objetividade de todos ou da maioria dos enunciados das questões foi percebida por 82,0 \%. 4.3.2 Componente de Conhecimento Específico Em relação aos enunciados das questões do Componente de Conhecimento Específico da prova, para 74,1 \% dos estudantes de Licenciatura avaliados da Área de Ciências Biológicas, a clareza e a objetividade ( Questão 5 ) estavam presentes em todas ou na maioria das questões ( Gráfico 4.9a, Gráfico 4.10a, e no Anexo II, a Tabela II.5a ).
A maioria dos estudantes de Licenciatura de todas as Grandes Regiões brasileiras considerou claros e objetivos todos ou a maioria dos enunciados das questões do Componente de Conhecimento Específico da prova, percentual sempre maior ou igual a 72,2 \%.
A diferença entre as regiões Nordeste e Sudeste é estatisticamente significativa.
Em relação aos enunciados das questões do Componente de Conhecimento Específico da prova, para 79,8 \% dos estudantes de Bacharelado avaliados da Área de Ciências Biológicas, a clareza e a objetividade ( Questão 5 ) estavam presentes em todas ou na maioria das questões ( Gráfico 4.9b, Gráfico 4.10b, e no Anexo II, a Tabela II.5b ).
A maioria dos estudantes de Bacharelado de todas as Grandes Regiões brasileiras considerou claros e objetivos todos ou a maioria dos enunciados das questões do Componente de Conhecimento Específico da prova, percentual sempre maior ou igual a 78,8 \%.
Não há diferença estatisticamente significativa entre as regiões.
A proporção de estudantes de Licenciatura que consideraram os enunciados das questões claros e objetivos apresenta uma tendência crescente em relação ao aumento de desempenho : mais elevada no quarto superior ( 80,2 \% ), se comparada ao quarto inferior de desempenho ( 66,7 \% ).
As diferenças entre todos os quartos de desempenho são estatisticamente significativas.
A proporção de estudantes de Bacharelado que consideraram os enunciados das questões claros e objetivos apresenta uma tendência crescente em relação ao aumento de desempenho : mais elevada no quarto superior ( 83,0 \% ), se comparada ao quarto inferior de desempenho ( 75,0 \% ).
Há diferenças estatisticamente significativas entre as proporções dos dois quartos inferiores de desempenho em relação aos dois superiores. 4.4 SUFICIÊNCIA DAS INFORMAÇÕES/INSTRUÇÕES FORNECIDAS Ao avaliarem as informações/instruções fornecidas para a resolução das questões ( Questão 6 ), 81,1 \% dos respondentes de Licenciatura da Área de Ciências Biológicas de todo o Brasil afirmaram que estas eram até excessivas ou suficientes em todas ou na maioria das questões ( Gráfico 4.11a, Gráfico 4.12a, e, no Anexo II, a Tabela II.6a ).
Quanto à distribuição de respondentes pelas Grandes Regiões, observa -se que a proporção de estudantes que consideraram as informações/instruções fornecidas até excessivas ou suficientes em todas ou na maioria das questões foi sempre superior ou igual a 77,4 \%, chegando a 84,0 \% na região Sudeste.
A diferença entre o maior percentual, obtido na região Sudeste, e o menor, da região Nordeste ( 77,4 \% ), é estatisticamente significativa, bem como as diferenças entre as regiões Norte e Sudeste e da região Nordeste em relação as regiões Sul e Centro-Oeste.
Ao avaliarem as informações/instruções fornecidas para a resolução das questões ( Questão 6 ), 85,7 \% dos respondentes de Bacharelado da Área de Ciências Biológicas de todo o Brasil afirmaram que estas eram até excessivas ou suficientes em todas ou na maioria das questões ( Gráfico 4.11b, Gráfico 4.12b, e, no Anexo II, a Tabela II.6b ).
Quanto à distribuição de respondentes pelas Grandes Regiões, observa -se que a proporção de estudantes que consideraram as informações/instruções fornecidas até excessivas ou suficientes em todas ou na maioria das questões foi sempre superior ou igual a 83,5 \%, chegando a 87,6 \% na região Centro-Oeste.
Não há diferenças estatisticamente significativas entre as regiões.
Levando -se em conta o desempenho dos participantes de Licenciatura, notam -se diferenças estatisticamente significativas entre as opiniões de estudantes de Licenciatura de todos os quartos de desempenho, como mostra o Gráfico 4.12a.
O percentual de participantes que avaliaram as informações/instruções como até excessivas ou suficientes em todas ou na maioria das questões foi mais elevado no quarto superior ( 87,0 \% ), percentual superior à média nacional ( 81,1 \% ).
Já no quarto inferior, a suficiência das informações/instruções declarada como até excessiva, em todas ou na maioria das questões foi percebida por 73,9 \% dos respondentes.
Essas proporções são crescentes com os quartos de desempenho.
Levando -se em conta o desempenho dos participantes de Bacharelado, notam -se diferenças estatisticamente significativas entre as opiniões de estudantes de Bacharelado dos dois quartos inferiores em relação aos dois quartos superiores, como mostra o Gráfico 4.12b.
O percentual de participantes que avaliaram as informações/instruções como até excessivas ou suficientes em todas ou na maioria das questões foi mais elevado no quarto superior ( 88,6 \% ), percentual superior à média nacional ( 85,7 \% ).
Já no quarto inferior, a suficiência das informações/instruções declarada como até excessiva, em todas ou na maioria das questões foi percebida por 81,7 \% dos respondentes.
Essas proporções são crescentes com os quartos de desempenho. 4.5 DIFICULDADE ENCONTRADA AO RESPONDER À PROVA \textbf{Indagados} sobre as dificuldades com as quais se depararam ao responderem à prova ( Questão 7 ), 18,2 \% dos estudantes de Licenciatura apontaram o Desconhecimento do conteúdo.
Para 49,5 \%, a Forma diferente de abordagem do conteúdo foi indicada como dificuldade.
Já a Falta de motivação para fazer a prova foi a dificuldade apontada por 16,6 \% dos respondentes.
Considerando -se todo o Brasil, 11,5 \% dos respondentes afirmaram que não tiveram qualquer tipo de dificuldade para responder à prova ( Tabela II.7a no Anexo II ).
Os Gráficos 4.13a e 4.14a apresentam os percentuais de estudantes de Licenciatura que apontaram o Desconhecimento do conteúdo como dificuldade percebida ao responderem à prova.
Na análise por Grandes Regiões, o percentual de inscritos e presentes de Licenciatura que apontaram o Desconhecimento do conteúdo como dificuldade ao responder à prova não superou 20,8 \%.
Os percentuais variaram de 17,2 \% na região Sudeste a 20,8 \% na Norte, sendo essa diferença significativa estatisticamente.
A Forma diferente de abordagem do conteúdo foi a escolha modal dos estudantes de Licenciatura, com percentuais que variaram de 47,2 \% ( região Sudeste ) a 51,7 \% ( Nordeste ).
O percentual de alunos que citaram a Falta de motivação como dificuldade variou de 11,7 \% ( região Norte ) a 18,9 \% ( região Sudeste ).
Os que declararam não ter qualquer tipo de dificuldade para responder à prova variaram de 10,4 \%, nas regiões Norte e Nordeste, a 12,8 \%, nas regiões Sudeste e Centro-Oeste. \textbf{Indagados} sobre as dificuldades com as quais se depararam ao responderem à prova ( Questão 7 ), 19,9 \% dos estudantes de Bacharelado apontaram o Desconhecimento do conteúdo.
Para 45,0 \%, a Forma diferente de abordagem do conteúdo foi indicada como dificuldade.
Já a Falta de motivação para fazer a prova foi a dificuldade apontada por 18,6 \% dos respondentes.
Considerando -se todo o Brasil, 12,8 \% dos respondentes afirmaram que não tiveram qualquer tipo de dificuldade para responder à prova ( Tabela II.7b no Anexo II ).
Os Gráficos 4.13b e 4.14b apresentam os percentuais de estudantes de Bacharelado que apontaram o Desconhecimento do conteúdo como dificuldade percebida ao responderem à prova.
Na análise por Grandes Regiões, o percentual de inscritos e presentes de Bacharelado que apontaram o Desconhecimento do conteúdo como dificuldade ao responderem à prova não superou 21,5 \%.
Os percentuais variaram de 15,0 \% na região Norte a 21,5 \% na Nordeste, não havendo diferenças significativa estatisticamente entre as regiões.
A Forma diferente de abordagem do conteúdo foi a escolha modal dos estudantes de Bacharelado, com percentuais que variaram de 42,9 \% ( região Nordeste ) a 56,2 \% ( Norte ).
O percentual de alunos que citaram a Falta de motivação como dificuldade variou de 13,9 \% ( região Norte ) a 19,0 \% ( região Sudeste ).
Os que declararam não ter qualquer tipo de dificuldade para responder à prova variaram de 12,3 \% na região Nordeste a 13,1 \% na Sudeste.
Com relação aos quartos de desempenho, o Desconhecimento do conteúdo foi a opção escolhida por 18,9 \% dos estudantes de Licenciatura do quarto superior e 16,7 \% do quarto inferior.
Não há diferenças estatisticamente significativas entre os quartos de desempenho.
A alternativa modal para os alunos de Licenciatura, quando agregados pelos quartos de desempenho, foi que a dificuldade encontrada foi causada pela Forma diferente de abordagem do conteúdo : 50,2 \% no quarto inferior e 44,7 \% do quarto superior escolheram essa alternativa.
Com relação aos quartos de desempenho, o Desconhecimento do conteúdo foi a opção escolhida por 21,5 \% dos estudantes de Bacharelado do quarto superior e 18,8 \% do quarto inferior.
Não há diferenças estatisticamente significativas entre os quartos de desempenho.
A alternativa modal para os alunos de Bacharelado, quando agregados pelos quartos de desempenho, foi que a dificuldade encontrada foi causada pela Forma diferente de abordagem do conteúdo : 46,3 \% no quarto inferior e 41,3 \% do quarto superior escolheram essa alternativa. 4.6 CONTEÚDOS DAS QUESTÕES OBJETIVAS DA PROVA Ao se analisarem os conteúdos das questões objetivas da prova ( Questão 8 ), verifica -se que um percentual muito pequeno dos estudantes de Licenciaturas avaliados, apenas 5,2 \%, afirmou que não estudou ainda a maioria desses conteúdos ( Gráficos 4.15a, Gráfico 4.16a, e a Tabelas II.8a no Anexo II ).
A maioria ( 62,4 \% ) afirmou ter estudado e aprendido muitos ou todos os conteúdos avaliados.
Na análise por Grande Região, a proporção de respondentes de Licenciatura que escolheram a opção não estudou ainda a maioria desses conteúdos foi pequena.
Observa -se que nas regiões Norte ( 6,7 \% ), Nordeste ( 6,6 \% ) e Sul ( 5,9 \% ), apesar de pequenas, as proporções foram maiores do que a média nacional ( 5,2 \% ).
Observa -se diferença estatisticamente significativa entre a região Sudeste e as regiões Norte, Nordeste e Sul.
Em todas as regiões, a maioria dos presentes afirmou ter estudado e aprendido muitos ou todos os conteúdos, com proporções variando entre 54,9 \% na região Norte e 70,1 \% na Sudeste.
Ao se analisarem os conteúdos das questões objetivas da prova ( Questão 8 ), verificou -se que um percentual muito pequeno dos estudantes de Bacharelados avaliados, apenas 2,5 \%, afirmou que não estudou ainda a maioria desses conteúdos ( Gráficos 4.15b, Gráfico 4.16b, e a Tabelas II.8b no Anexo II ).
A maioria ( 71,8 \% ) afirmou ter estudado e aprendido muitos ou todos os conteúdos avaliados.
Na análise por Grande Região, a proporção de respondentes de Bacharelado que escolheram a opção não estudou ainda a maioria desses conteúdos foi pequena.
Observa -se que, com exceção da região Sudeste ( 2,2 \% ), nas demais regiões, apesar de pequenas, as proporções foram maiores do que a média nacional ( 2,5 \% ).
Não há diferenças estatisticamente significativas entre as regiões.
Em todas as regiões, a maioria dos presentes afirmou ter estudado e aprendido muitos ou todos os conteúdos, com proporções variando entre 68,2 \% na região Nordeste e 73,8 \% na Sudeste.
Considerando -se separadamente as opiniões de estudantes de Licenciatura dos quatro quartos de desempenho, observa -se que, no quarto inferior, 10,3 \% ofereceram como resposta que não estudou ainda a maioria desses conteúdos, sendo 1,3 \% os do quarto superior com a mesma resposta.
A diferença entre os alunos que optaram por esse motivo de dificuldade entre todos os quartos de desempenho é estatisticamente significativa.
Tendo em conta o quarto superior, 77,7 \% dos alunos de Licenciatura afirmaram ter estudado e aprendido muitos ou todos os conteúdos.
No outro extremo, no primeiro quarto, 46,1 \% optaram pelas mesmas categorias.
As proporções são crescentes com o desempenho.
Considerando -se separadamente as opiniões de estudantes de Bacharelado dos quatro quartos de desempenho, observa -se que, no quarto inferior, 4,8 \% ofereceram como resposta que não estudou ainda a maioria desses conteúdos, sendo 1,4 \% os do quarto superior com a mesma resposta.
A diferença entre os alunos que optaram por esse motivo de dificuldade entre o primeiro quarto de desempenho e os demais é estatisticamente significativa.
Tendo em conta o quarto superior, 82,6 \% dos alunos de Bacharelado afirmaram ter estudado e aprendido muitos ou todos os conteúdos.
No outro extremo, no primeiro quarto, 60,5 \% optaram pelas mesmas categorias.
As proporções são crescentes com o desempenho. 4.7 TEMPO GASTO PARA CONCLUIR A PROVA Ao responderem sobre o tempo de conclusão da prova ( Questão 9 ), mais de três quartos dos estudantes de Licenciatura ( 77,2 \% ) afirmaram ter gasto Entre duas e quatro horas ( Gráfico 4.17a, Gráfico 4.18a e, no Anexo II, a Tabela II.9a ).
Considerando -se as cinco Grandes Regiões brasileiras, o percentual dos que utilizaram Entre duas e quatro horas para finalizar a prova foi inferior ao percentual nacional em quatro delas, sendo superior apenas na região Sudeste ( 78,4 \% ).
Nas Grandes Regiões, o percentual de alunos de Licenciatura que \textbf{dispensaram} Entre duas e quatro horas para concluir a prova ficou igual ou acima de 75,6 \%, como mostra o Gráfico 4.17a.
Não há diferenças estatisticamente significativas entre as proporções das Grandes Regiões.
Ao responderem sobre o tempo de conclusão da prova ( Questão 9 ), mais de três quartos dos estudantes de Bacharelado ( 77,6 \% ) afirmaram ter gasto Entre duas e quatro horas ( Gráfico 4.17b, Gráfico 4.18b e, no Anexo II, a Tabela II.9b ).
Considerando -se as cinco Grandes Regiões brasileiras, os percentuais dos que utilizaram Entre duas e quatro horas para finalizar a prova nas regiões Nordeste ( 76,4 \% ) e Sul ( 75,2 \% ) foram inferiores ao percentual nacional.
Nas demais Grandes Regiões, o percentual de alunos de Bacharelado que \textbf{dispensaram} Entre duas e quatro horas para concluir a prova ficou igual ou acima de 78,0 \%, como mostra o Gráfico 4.17b.
Há diferenças estatisticamente significativas entre a proporção do Sul e Centro-Oeste.
Uma vez sendo analisadas as alternativas escolhidas pelos estudantes de Licenciatura que se situam nos diferentes quartos de desempenho, observa -se uma tendência crescente : uma proporção maior de participantes no quarto superior declarou ter gasto Entre duas e quatro horas para concluir a prova quando comparadas com os dos quartos inferiores.
Há diferenças estatisticamente significativas entre todos os quartos de desempenho.
Uma vez sendo analisadas as alternativas escolhidas pelos estudantes de Bacharelado que se situam nos diferentes quartos de desempenho, observa -se uma tendência crescente : uma maior proporção de participantes no quarto superior declarou ter gasto Entre duas e quatro horas para concluir a prova quando comparadas com os dos quartos inferiores.
As diferenças entre o primeiro ( 68,9 \% ) e os demais quartos de desempenho são estatisticamente significativas, bem como entre o segundo e o quarto superior de desempenho.

% matched lemmas: arredondamento, crescer, dispensar, indagar, numérico, ordenar, percentagem, questionário, recorte, reprodução, subconjunto
\end{textsample}
